\chapter{Binary Lifting and Cross-Architecture Translation}
\label{ch:binary-lifting}

A linker produces bytes that a processor can execute, but a binary optimizer needs more than bytes.
It needs explicit definitions, uses, control edges, memory effects, and instruction semantics.
\term{Binary lifting} translates machine instructions into an IR suitable for analysis or further translation.
It recovers an executable description of machine behavior, not the original source program.
Names, types, loop structure, and source abstractions may be absent or ambiguous after compilation.

The small lifter in this chapter handles three real x86 byte instructions.
Its narrow scope makes every state update inspectable and enables comparison against hardware.
The surrounding architecture discussion explains what a general lifter would additionally need for control flow, memory, exceptions, and translation to another ISA.

\section{Choose the observable machine state}
\label{sec:ch64-state}

For the example, state contains a 64-bit \code{RAX} value and six arithmetic flags.
Each instruction changes the low byte \code{AL}; the upper $56$ bits stay unchanged.
The supported byte sequences are listed in Table~\ref{tab:ch64-subset}.
Every instruction has length two, and unsupported opcodes or truncated immediates are rejected.
Prefixes, memory operands, branches, and instruction-fetch faults are outside this executable subset.

\begin{table}[htbp]
\caption{The lifter accepts three fixed-form byte operations with distinct flag contracts.}
\label{tab:ch64-subset}
\centering
\begin{tabular}{lll}
\toprule
Bytes & Operation & Arithmetic flags \\
\midrule
\code{B0 ib} & move immediate to AL & unchanged \\
\code{04 ib} & add immediate to AL & all six computed \\
\code{34 ib} & exclusive OR immediate with AL & AF unspecified; others computed \\
\bottomrule
\end{tabular}
\end{table}

The encoding and flag distinctions are documented in the instruction reference~\cite{cloutier-2026-x86}.
In particular, XOR clears carry and overflow, computes sign, zero, and parity, and leaves auxiliary carry undefined.
The model represents that last condition by an unspecified marker rather than choosing a fake architectural value.
A later operation that observes an unspecified flag would need nondeterministic semantics or an explicit restriction.
None of our accepted instructions reads that flag.

For byte addition, let $a,b\in\{0,\ldots,255\}$ and $r=(a+b)\bmod256$.
Carry is true when $a+b\ge256$; zero is true when $r=0$; sign is bit seven of $r$; parity is true when $r$ has an even number of set bits.
Auxiliary carry records carry out of bit three.
Signed overflow is true when the operands have the same sign bit and the result has the other sign bit.
These predicates distinguish $255+1$ from $127+1$: the former has unsigned carry without signed overflow, and the latter has signed overflow without unsigned carry.

\section{Lift writes into single-assignment values}
\label{sec:ch64-lift}

The lifter gives every new low-byte value a distinct name.
It also retains the source byte offset for provenance.
Algorithm~\ref{alg:ch64-lift} describes the straight-line conversion; flag updates are interpreted according to the lifted opcode's contract.
For $n$ accepted instructions, decoding and row construction take $O(n)$ time and storage.
Unlike a general x86 decoder, this one never guesses an instruction boundary.

\begin{algorithm}[htbp]
\caption{Lift the fixed two-byte AL instruction subset.}
\label{alg:ch64-lift}
\begin{pseudocode}
\Require byte buffer, known entry at offset zero, accepted opcode profile
\Ensure SSA value rows with original offsets, or explicit rejection
\State $pc\gets0$; $current\gets$ entry low-byte value
\While{$pc$ is below the buffer length}
  \State reject unless opcode is accepted and one immediate byte remains
  \State allocate fresh value name $v$
  \State emit the opcode's byte operation with source $current$ and immediate
  \State associate its flag contract and source offset $pc$
  \State $current\gets v$; $pc\gets pc+2$
\EndWhile
\State return the lifted rows and final low-byte value
\end{pseudocode}
\end{algorithm}

Start with \code{RAX=0x12345600} and execute bytes \code{B0 FF 04 01}.
The first operation produces low byte $255$ without changing flags.
The second produces low byte zero, carry one, zero one, sign zero, parity one, auxiliary carry one, and overflow zero.
The final full register is \code{0x12345600}, because neither byte write clears the upper bits.
Table~\ref{tab:ch64-flags} includes a second addition that exercises signed overflow.

\begin{table}[htbp]
\caption{Unsigned carry and signed overflow describe different properties of byte addition.}
\label{tab:ch64-flags}
\centering
\begin{tabular}{rrrrrrrrr}
\toprule
$a$ & $b$ & $r$ & CF & PF & AF & ZF & SF & OF \\
\midrule
$255$ & $1$ & $0$ & $1$ & $1$ & $1$ & $1$ & $0$ & $0$ \\
$127$ & $1$ & $128$ & $0$ & $0$ & $1$ & $0$ & $1$ & $1$ \\
$128$ & $128$ & $0$ & $1$ & $1$ & $0$ & $1$ & $0$ & $1$ \\
\bottomrule
\end{tabular}
\end{table}

Figure~\ref{fig:ch64-register} shows the register update as a merge of unchanged high bits and a new low byte.
A translator that treats every subregister write as a full-width assignment would lose the high bits.
A translator that generalizes this byte rule to every x86 register width would also be wrong; width-specific behavior belongs in the machine semantics.

\begin{figure}[htbp]
\centering
\begin{tikzpicture}[node distance=14mm and 16mm]
\node[basic block] (high) {old RAX bits 63..8};
\node[basic block,right=of high] (low) {new AL bits 7..0};
\node[basic block,below=of high] (merge) {merge unchanged high bits};
\node[basic block,below=of merge] (out) {new full RAX};
\draw[cfg edge] (high)--(merge); \draw[cfg edge] (low)--(merge); \draw[cfg edge] (merge)--(out);
\end{tikzpicture}
\caption{The lifted byte operation preserves register state outside its architectural write footprint.}
\label{fig:ch64-register}
\end{figure}

\section{Recovered control flow has an evidence boundary}
\label{sec:ch64-cfg}

General binaries contain direct branches, indirect branches, calls, returns, embedded data, and sometimes code that changes after loading.
A decoder can identify the instruction at a supplied address without proving that the address is reachable code.
A control-flow recovery process must distinguish discovered successors from unknown successors.
Dropping an unresolved indirect edge is not a conservative analysis.

Algorithm~\ref{alg:ch64-discover} sketches a worklist for a fixed code image.
It reports unresolved targets explicitly and assumes any indirect-target analysis it uses is sound for its stated environment.
With a finite set of $N$ discovered instruction starts and $E$ supplied successor edges, graph bookkeeping is $O(N+E)$, excluding decoding, target analysis, and overlap handling.
For variable-length ISAs, two entries can lead to overlapping instruction interpretations; a representation may preserve those interpretations or reject that input profile, but must not silently merge incompatible decodes.

\begin{algorithm}[htbp]
\caption{Discover control flow while retaining unresolved behavior.}
\label{alg:ch64-discover}
\begin{pseudocode}
\Require immutable code image, trusted entry set, instruction semantics
\Ensure discovered instruction graph and explicit unresolved exits
\State initialize worklist with entries and an empty visited set
\While{the worklist is nonempty}
  \State remove address $a$; skip it if already visited
  \State decode and lift at $a$, or record a decoding failure
  \State record direct successors and soundly established indirect successors
  \If{some possible successor remains unresolved}
    \State retain an unknown exit with the required machine-state contract
  \EndIf
  \State enqueue unvisited known successors
\EndWhile
\end{pseudocode}
\end{algorithm}

Dynamic translation can resolve an indirect destination at runtime and dispatch to a translated block.
Its cache key must include everything that affects decoding and semantics, such as code identity and execution mode.
Self-modifying code requires invalidation or a stronger immutable-code assumption.
Fault attribution and asynchronous observation may require mappings from translated instructions back to guest instruction boundaries.

\section{Translate semantics, then lower to a host}
\label{sec:ch64-cross-isa}

A \term{cross-architecture translator} composes a guest decoder and semantic lifter with a host backend and runtime services.
Figure~\ref{fig:ch64-bridge} makes the intermediate machine state explicit.
The host need not have the same flags or register widths; it may compute guest flags as ordinary values and omit computations proved dead.
The proof obligation is preservation of guest-visible behavior, not one host instruction per guest instruction.

\begin{figure}[htbp]
\centering
\begin{tikzpicture}[node distance=13mm]
\node[basic block] (guest) {guest bytes and machine state};
\node[basic block,below=of guest] (ir) {explicit semantic IR};
\node[basic block,below=of ir] (host) {host instructions and runtime};
\draw[cfg edge] (guest)--node[right,font=\small] {lift}(ir);
\draw[cfg edge] (ir)--node[right,font=\small] {lower}(host);
\end{tikzpicture}
\caption{Cross-ISA translation composes two semantic boundaries and an environment model.}
\label{fig:ch64-bridge}
\end{figure}

Memory ordering, endianness, alignment, atomicity, system calls, and ABI details can break a translator even when integer expressions are correct.
An emulator may provide an entire guest environment, while a process translator may adapt selected operating-system interfaces.
Neither scope follows automatically from choosing an IR.
Binary optimization likewise needs to preserve live registers and flags at exits, even if no recovered source variable appears to use them.

\section{Verification and exercises}
\label{sec:ch64-exercises}

The hardware oracle in \path{example-code/architecture-labs/flags.c} runs byte addition for all $65{,}536$ operand pairs and captures the six arithmetic flags.
The Python tests compare those observations with the lifted evaluator, including preservation of high register bits.
The oracle uses register-register addition to exercise the same arithmetic semantics; encoding acceptance is a separate decoder check.
GCC and Clang builds, with and without address/undefined-behavior sanitizers, provide four configurations.
This experiment does not validate a general decoder, a memory lifter, or a host code generator.

\begin{enumerate}
\item\label{ex:ch64-add} Compute all six modeled flags for $128+128$ and explain why both carry and overflow are set.
\item\label{ex:ch64-partial} Starting with \code{RAX=0xABCD12}, execute \code{B0 07 04 09}.
Give the final register and explain why the upper bytes survive.
\item\label{ex:ch64-unknown} An indirect branch has two observed targets but no proof that they are exhaustive.
Explain why a static optimizer cannot use those observations alone to delete the unknown exit.
\item\label{ex:ch64-xor} Which flag cannot be compared to a fixed expected bit after XOR under the documented contract?
Explain the difference between an unspecified result and preserving the previous value.
\end{enumerate}

\section{Further reading}
\label{sec:ch64-reading}

BAP provides a published architecture for binary analysis through lifted representations~\cite{brumley-2011-bap}.
The instruction reference gives the encodings and state effects used in the small experiment~\cite{cloutier-2026-x86}.
